\documentclass[10pt,a4paper]{amsart}
\usepackage[margin=19mm]{geometry}
\usepackage{amsmath,amssymb,mathtools}
\usepackage[scr=rsfs,cal=euler]{mathalfa}
\usepackage{xfrac}
\usepackage[unicode,hidelinks,bookmarks=false]{hyperref}
\newcommand{\ie}{i.e.\ }
\newcommand{\cf}{cf.\ }
\newcommand{\resp}{respectively\ }
\newcommand{\Subsection}{\subsection}
\providecommand{\nobreakdash}{\nobreak\hbox{-}}
\setlength{\emergencystretch}{2em}
\multlinegap0pt
\usepackage[shortlabels]{enumitem}
\usepackage{amssymb}
\usepackage{tikz}
\newtheorem{proposition}{Proposition}
\def\KK{{\mathcal{K}}}
\def\T{{\mathbb{T}}}
\title{Optimal Dynamical Frames}
\author{A. Aguilera, C. Cabrelli, F. Negreira, V. Paternostro}
\date{Source: arXiv:2506.00567v3 (CC BY 4.0)}
\begin{document}
\maketitle

\noindent\emph{Source and licence.} Optimal Dynamical Frames. Version arXiv:2506.00567v3 (CC BY 4.0), distributed by arXiv under the Creative Commons Attribution 4.0 International licence, http://creativecommons.org/licenses/by/4.0/, as recorded in the arXiv licence metadata for this exact version and shown on the abstract page https://arxiv.org/abs/2506.00567v3. Full licence notice: \texttt{licenses/arxiv-2506.00567-CC-BY-4.0.txt}.\par\smallskip

\noindent\textit{Editorial scope.} Counted unit: the article's proposition dimcoro2 (source0.tex lines 1236--1237). The article's complete original proof of it is delivered with it (source0.tex lines 1238--1280, one \texttt{proof} environment of 43 lines). No mathematics is written, repaired or re-derived: the statement and the proof are the article's own lines, in the article's own notation, and no retained line is substituted or re-flowed.  Retained uncounted as the source-local context the proof needs: lines 1212--1217.  Nothing was omitted from any retained span, so the package carries no omission record.  No retained line contains a citation, so the wrapper carries no bibliography and no bibliography entry.  No retained line contains a figure environment, an image-inclusion command or drawing code, so no image is delivered and the package declares no asset dependency.  Editorial formatting only, in addition: an AMS \texttt{amsart} preamble that restates the article's own declarations and macros used by the retained lines, namely the environments \texttt{proposition} on separate counters, together with the standard packages those lines need.  The retained environments print in this document as Proposition 1 in order of appearance, whereas the article prints the same items with the numbering of its own full document.  The complete native source of the article as distributed in the arXiv e-print for this exact version is bundled unchanged as \texttt{sources/179093/source0.tex}.
\begin{proposition}\label{dimcoro}
For the Parseval index of $S^*|_N$ and $A_N$ we have that
\begin{equation*}
\gamma_p(S^*|_N)=\gamma_p(A_N)=\dim\mathcal{L}.
\end{equation*}
\end{proposition}

\begin{proposition}\label{dimcoro2}
Under the standing setting of this section, we have $\mathcal{L} = W_1$.\end{proposition}

\begin{proof}
From Proposition \ref{dimcoro} we know that $\gamma_p(S^*|_N)=\dim\overline{(U-A_N)(N)}$. Note that
\begin{equation*}
\overline{(U-A_N)(N)}=\overline{\text{rg}(U-P_NU)P_N}
\end{equation*}
where the operator $(U-P_NU)P_N$ is understood as acting in $L^2(\T,\KK)$. Noting that $(U-P_NU)P_N=(I-P_N)UP_N=P_{N^\perp} UP_N$ where $N^\perp=L^2(\T,\KK)\ominus N$ and $P_{N^\perp}$ is the projection in $L^2(\T,\KK)$ onto $N^\perp$, we have that
\begin{equation*}
\overline{\text{rg}(U-P_NU)P_N}=\ker(P_NU^*P_{N^\perp})^\perp.
\end{equation*}

Let us now study the kernel of $P_NU^*P_{N^\perp}:L^2(\T,\KK)\to L^2(\T,\KK)$. To that end, we decompose $L^2(\T,\KK)$ as follows
\begin{equation*}
L^2(\T,\KK)=\left(L^2(\T,\KK)\ominus H^2_\KK\right)\oplus N\oplus S(M)\oplus W
\end{equation*}
where, recall, $M=H^2_\KK\ominus N$ and $W=M\ominus S(M)$. 

First, note that since $L^2(\T,\KK)\ominus H^2_\KK$ is a $U^*$-invariant subspace of $N^\perp$ then
\begin{equation*}
U^*P_{N^\perp}(L^2(\T,\KK)\ominus H^2_\KK)=U^*(L^2(\T,\KK)\ominus H^2_\KK)\subseteq L^2(\T,\KK)\ominus H^2_\KK\subseteq N^\perp,
\end{equation*}
and thus $P_NU^*P_{N^\perp}(L^2(\T,\KK)\ominus H^2_\KK))=\{0\}$. 

Second, the inclusion $N\subseteq\ker(P_NU^*P_{N^\perp})$ follows from $P_{N^\perp}(N)=\{0\}$. 

Third, note that as $S(M)=U(M)\subseteq N^\perp$ then
\begin{equation*}
U^*P_{N^\perp}(S(M)) =U^*U(M)=M,
\end{equation*}
and thus $P_NU^*P_{N^\perp}(S(M))=\{0\}$.

Finally, let $f\in W$ and note that
\begin{equation*}
P_NU^*P_{N^\perp} f=P_NU^*f=P_NS^*f=S^*f
\end{equation*} 
because $W\subseteq N^\perp$, $S^*(W)\subseteq N$ and $N$ is $S^*$-invariant. Thus, if we decompose $W=\KK_0\oplus W_1$, where as before $\KK_0=W\cap\KK$, and we recall that $\ker S^*=\KK$, we have, altogether,
\begin{equation*}
\ker(P_NU^*P_{N^\perp})=\left(L^2(\T,\KK)\ominus H^2_\KK\right)\oplus N\oplus S(M)\oplus \KK_0.
\end{equation*}



It follows that $\ker(P_NU^*P_{N^\perp})^\perp=W_1$ and the proof is finished.
\end{proof}

\end{document}
